\begin{proposition}[B.4.9 (Proved): ordinary logarithmic torus cocycles and the remaining descent classes]
Let $\mathcal T\to U$, $U=\Delta_q^*\times\Delta_r^*$,
be the Jacobian two-torus family of the marked cusp,
after $w=r^6$.  A holomorphic homomorphism of group
families
\[
 \C^*\times U\longrightarrow\mathcal T
\]
is determined by a flat integral cycle.  For this family
all such cycles are
\[
 c=a\widehat w+b\widehat\delta,\qquad a,b\in\Z.
\]
In the exponential cusp coordinates its homomorphism
is $z\mapsto(z^a,z^b)$.  Thus the ordinary Cartier
logarithm construction along a globally specified
cocharacter has transition maps
\[
 (Z,W)\longmapsto(g_{ij}^{\,a}Z,g_{ij}^{\,b}W).
\]
If these units extend over the prescribed fan charts,
the maps extend there and preserve each toric
boundary component.  Changing only these transition
units cannot repair the component permutations
$+4,-3,0$ refuted in B.3.1.

This is not a classification of all torus-torsor
descent.  In fact, for the sheaf $\mathscr T$ of
holomorphic sections of $\mathcal T$,
\[
 H^1(U,\mathscr T)\simeq H^2(U,\Lambda)\simeq\Z^2.
\]
The last group is represented by the coinvariant
classes of $\widehat\gamma,\widehat u$.  The ordinary
fixed-cocharacter Cartier constructions have zero
image in this group.  The actual cusp family with
the coherent closing translations of C.3.5 also
has zero torsor class, by the explicit holomorphic
section below.  Thus nonzero classes in this
$\Z^2$ do not have its prescribed marked punctured
cusp germ.  No compactification of an arbitrary
torsor is asserted.
\end{proposition}
\begin{proof}
First consider any smooth holomorphic family of
complex tori with integral period local system
$\Lambda$ and relative Lie bundle $\mathcal V$.
Lift a group homomorphism from $\C^*$ along the
universal covering $z=\exp(2\pi i t)$.
The lifted homomorphism from the additive group
$\C$ to the universal covering vector space of
each torus is $t\mapsto t v$ for a holomorphic
section $v$ of $\mathcal V$.  Periodicity under
$t\mapsto t+1$ says precisely that $v$ belongs
to the period lattice in every fibre.  In any
local marking its integral coefficients are
continuous, and hence locally constant.  It
therefore defines a flat global section of
$\Lambda$.  Conversely any such section gives
the homomorphism by this formula.  The image
need not be a closed elliptic subgroup; this
argument applies without a polarization and
also to simple two-tori.

In the ordered basis
$(\widehat\gamma,\widehat u,\widehat w,\widehat\delta)$,
the actual commuting linear monodromies have
nilpotent parts
\[
 N_q=A_0-I=
 \begin{pmatrix}
 0&0&0&0\\0&0&0&0\\0&1&0&0\\-1&0&0&0
 \end{pmatrix},\qquad
 N_r=-12\widehat\delta\otimes\gamma.
\]
The sign of $N_r$ is the positive-loop closing
convention of C.3.5; reversing it leaves all
groups below unchanged.  Their common integral
kernel is exactly $\Z\widehat w\oplus\Z\widehat\delta$.
The two corresponding columns of $\Pi_c$ are
$(1,0)^t,(0,1)^t$.  Exponentiating the homomorphism
therefore gives $(z^a,z^b)$, as asserted.

Fujimoto's logarithmic cocycle formula
\citep[Definition 1.1]{Fujimoto1990AB} works in
the same way for a complex torus of any dimension:
apply the displayed homomorphism to the transition
functions $g_{ij}$ of a line bundle.  Equivalently
use $c\log(g_{ij})/(2\pi i)$ modulo the period
lattice.  A different logarithm adds an integral
cycle, and the triple-overlap defect is an
integral cycle as well.  With the present common
invariant cycles this is exactly multiplication
by $(g_{ij}^a,g_{ij}^b)$.  Torus multiplication
extends to every toric variety of a fixed fan
and preserves every torus orbit and boundary
ray.  For an extending unit it is an automorphism
on the full chart.  Consequently it contributes
no permutation of the central component labels.
The product-of-permutations obstruction of B.3.1
is unchanged.  This conclusion keeps the fan,
period action, and local component models fixed.

To exhibit the scope left outside that calculation,
retract $U$ to its real two-torus.  The cellular
cochain complex with coefficients in the local
system is the two-generator Koszul complex.
Its last differential has image
$N_q\Lambda+N_r\Lambda$, and hence
\[
 H^2(U,\Lambda)
   =\Lambda/(N_q\Lambda+N_r\Lambda)
   =\Lambda/(\Z\widehat w+\Z\widehat\delta)
   \simeq\Z\widehat\gamma\oplus\Z\widehat u.
\]
The relation lattice is primitive: $N_q$ already
has image $\Z\widehat w+\Z\widehat\delta$.
The exact matrix and Smith calculation is saved
in \path{code/ab_push_hopf_winding.py}.

The analytic exponential sequence is
\[
 0\longrightarrow\Lambda\longrightarrow
 \mathcal O(\mathcal V)\longrightarrow
 \mathscr T\longrightarrow0.
\]
The product of two punctured discs is Stein,
and $\mathcal O(\mathcal V)$ is coherent.
Cartan's theorem B therefore kills its positive
cohomology
\citep[Section 5]{Cartan1952B_AB}.
The long exact sequence gives the stated
isomorphism $H^1(U,\mathscr T)\simeq H^2(U,\Lambda)$.
These classes do define holomorphic torsors by
translation gluing, but this calculation contains
no compactness assertion at either missing axis.

For the last assertion, naturality of the two
exponential sequences sends the first Chern class
of a Cartier line bundle to its multiple of the
flat cycle $c$ in $H^2(U,\Lambda)$.  Both common
invariant cycles lie in $N_q\Lambda$, so this
multiple is zero.  A holomorphic trivializing
gauge over $U$ need not extend at the missing
axes.  The proof thus neither identifies
arbitrary descent data with monomial cocycles
nor asserts a meromorphic comparison of
arbitrary analytic fillings.

We finally compute the actual affine class rather
than infer it from these linear monodromies.  Retain
the cusp functions of C.3.5,
\[
 q=e^{2\pi i s},\qquad
 C_0(q)=\beta^{\rm part}+s,\qquad
 \ell=-s/12+\kappa_0(q),
 \qquad D(q)=C_0(q)-12\kappa_0(q).
\]
The functions $\mu,C_0,\kappa_0$ are single-valued
and holomorphic at $q=0$.  The actual cusp unit
representative of the even-shift translation is
\[
 \bar b_{h,c}(q)=
 \left(-\frac h2\mu(q),\,
       -\frac h{12}(c+D(q))-\frac{h^2}{24}\right).
\]
On the logarithmic parameter cover define
\[
 \sigma(q,c)=
 \left(-\frac c2\mu(q),\,
       -\frac{c^2}{24}-\frac{cD(q)}{12}\right).
\]
Direct subtraction gives the exact identity
\[
 \sigma(q,c+h)-\sigma(q,c)=\bar b_{h,c}(q).
\]
In particular it is equivariant for the deck
transformation $c\mapsto c-12$ corresponding to
a positive $r$ loop, and it is single-valued in
$q$.  The original cusp has no extra affine
$q$ translation ($l_0=0$).  Thus $\sigma$ descends
to a holomorphic section of the actual affine
torus family.  Equivalently its $q$ and $r$
affine lifts commute exactly: the $r$ translation
$\bar b_{-12,c}$ is independent of the logarithm
branch $s$.  The resulting integral commutator
representing the Chern class is zero already in
$\Lambda$, before passage to coinvariants.

Multiplication by the exponential of $\sigma$
is a torus-unit automorphism of every cusp fan
chart for a fixed nonzero parameter.  Since its
coefficients are holomorphic in the full $q$
disc, this comparison includes the prescribed
Mumford cusp fibre for $r\ne0$.  The quadratic
dependence on $c=-(12/(2\pi i))\log r$ is retained:
it supplies no extension across $r=0$ and no
global comparison with the finite-place models.
Changing coherent closings by a scalar depending
only on the parameter has zero $q$ commutator
and leaves this Chern-class conclusion unchanged.
The explicit polynomial identity is checked in
\path{code/ab_push_hopf_winding.py}.
\end{proof}
